Dinic 每轮先用 BFS 建立分层图，再用 DFS 沿分层图增广。
数组 \texttt{cur} 记录当前弧，保证无效边不重复扫描。
建边用 \texttt{add(u,v,cap)}，最大流用 \texttt{flow(s,t)}。

最大流最小割定理：任意网络中，最大流值等于最小割容量。
因此把题目转成网络后，求最大流也就同时给出了最小割的值。

常见建图：

\begin{itemize}
  \item 二分图最大匹配：源点连左部，右部连汇点，中间边容量设为 $1$。
  \item 最小点割：拆点 $v_{in}\to v_{out}$，容量设为点权或 $1$。
  \item 多源汇：增加超级源、超级汇，连接原源和原汇。
\end{itemize}
